% =====================================================================
% introduccion.tex
% Responsabilidad única: presentar el laboratorio, sus preguntas, el
% requisito del curso y la contribución del proyecto. Sin resultados.
% =====================================================================
\section{Introducción}
\label{sec:introduccion}

Un laboratorio de microbiología recibe a diario muestras de clientes de
distintas líneas de negocio (cosmética, farmacéutica, alimentos, veterinaria y
centrales de mezclas), las siembra según el ensayo solicitado (mesófilos
aerobios, mohos y levaduras, coliformes, \textit{Salmonella}, esterilidad,
entre otros) y emite un informe que declara cada muestra conforme o no
conforme frente al límite de la norma. Cada muestra deja un rastro en la base
de datos del laboratorio, desde la programación de la visita de muestreo y las
condiciones de transporte hasta la lectura, el eventual reensayo y la fecha
del informe, de forma que esa base contiene información suficiente para
anticipar parte de lo que va a ocurrir en la operación.

\paragraph{Preguntas del laboratorio.} El trabajo se organiza alrededor de tres
preguntas operativas y una cuarta, descriptiva:

\begin{itemize}
    \item \textbf{Carga diaria:} ¿cuántas muestras llegarán el próximo día
          hábil?, de forma que se pueda planear el personal de recepción y
          siembra con algo mejor que el promedio de la semana anterior.
    \item \textbf{Riesgo de no conformidad:} ¿qué muestras, al momento de
          recibirlas, tienen mayor probabilidad de resultar no conformes?, con
          el objetivo de priorizarlas sin saturar al personal con alertas.
    \item \textbf{Falso positivo:} ¿qué tan probable es que un resultado no
          conforme sea falso (contaminación del ensayo y no del producto), y
          cuándo conviene reensayar?
    \item \textbf{Grupos:} ¿qué perfiles de muestras y qué segmentos de
          clientes existen?, como insumo descriptivo para la planeación y la
          relación con los clientes.
\end{itemize}

\paragraph{Requisito del curso.} El proyecto final de Aprendizaje Automático
de Máquina~I pide recorrer, sobre un caso propio, los temas del curso
\cite{muller2016}: preprocesamiento y análisis exploratorio, $k$ vecinos más
cercanos, regresión logística con curva ROC y elección de umbral, árboles de
decisión con poda por complejidad de costo, regresión lineal simple y
múltiple, regresores KNN y de árbol, validación cruzada con búsqueda en
rejilla, K-means, DBSCAN y el perceptrón multicapa. Todos los modelos se
construyen con scikit-learn \cite{pedregosa2011}, el sobremuestreo con
imbalanced-learn \cite{lemaitre2017} y el manejo de tablas con pandas
\cite{mckinney2010}.

\paragraph{Contribución.} Más que un modelo puntual, se entrega un
procedimiento reproducible y honesto de punta a punta, en el que primeramente
un contrato de datos declara en qué momento del proceso se conoce cada
columna (lo que impide fugas de información), progresivamente un cribado con
una regla escrita antes de ver resultados elige qué análisis merece estudiarse
a fondo, después un diseño de Taguchi reduce el costo de explorar
configuraciones sin perder las conclusiones del factorial completo, y
finalmente un protocolo temporal estricto (ajustar con entrenamiento, decidir
con validación y abrir la prueba una sola vez) evalúa los modelos congelados.
Cada cifra del informe sale de una macro que escriben los notebooks al
ejecutarse, de forma que el texto no puede quedar desactualizado frente a los
resultados.

El resto del documento sigue ese orden: la Sección~\ref{sec:datos} describe
los datos y su contrato, la Sección~\ref{sec:exploracion} la exploración, la
Sección~\ref{sec:metodologia} la metodología común, las
Secciones~\ref{sec:regresion} a~\ref{sec:agrupamiento} cada análisis, la
Sección~\ref{sec:evaluacion} la evaluación final en prueba y la
Sección~\ref{sec:conclusiones} las conclusiones y limitaciones.
